\documentclass[11pt]{article}
\usepackage[margin=1in]{geometry}
\usepackage{amsmath,amssymb,amsfonts,mathtools}
\usepackage{array,booktabs,enumitem,graphicx,xcolor,hyperref}
\usepackage[T1]{fontenc}
\usepackage[utf8]{inputenc}
\title{The Hausdorff dimension of the boundary of the Mandelbrot set and Julia
  sets}
\author{Source-backed arXiv sample}
\date{}
\begin{document}
\maketitle
\section{Figure Captions}

\begin{figure}[htbp]
\centering
\fbox{\rule{0pt}{1.15in}\rule{0.78\linewidth}{0pt}}
\caption{Figure 5\} Gluing these two curves of \$S\_0\textasciicircum{}-\$ (resp. \$S\_0\textasciicircum{}+\$), one obtains a topological cylinder \$\textbackslash{}Cal C\_0\textasciicircum{}-\$ (resp. \$\textbackslash{}Cal C\_0\textasciicircum{}+\$), called the \{\textbackslash{}it outgoing}
\label{fig:source-1}
\end{figure}


\begin{figure}[htbp]
\centering
\fbox{\rule{0pt}{1.15in}\rule{0.78\linewidth}{0pt}}
\caption{Figure 7\} Now let us consider the following sequence of maps: \$\$ \textbackslash{}aligned U\_i @> \textbackslash{} \textbackslash{}t\_0\textasciicircum{}\{-1}
\label{fig:source-2}
\end{figure}


\begin{figure}[htbp]
\centering
\fbox{\rule{0pt}{1.15in}\rule{0.78\linewidth}{0pt}}
\caption{Figure 8\} To analyze the bifurcation, we need to consider \$q\$ incoming and \$q\$ outgoing Ecalle cylinders \$\textbackslash{}Cal C\_0\textasciicircum{}\{k,+}
\label{fig:source-3}
\end{figure}


Figures~\ref{fig:source-1} through~\ref{fig:source-3} preserve source-backed captions while replacing external graphics with deterministic placeholders.
\end{document}
